\chapter{数值测试}

为了系统验证本文所采用数值方法在强间断、复杂波系及长时演化条件下的准确性、稳定性与鲁棒性，本章围绕若干具有代表性的标准算例开展数值实验，并以可重复的参数设置和统一的诊断指标对结果进行评估。测试内容覆盖相对论流体与磁流体中的关键物理过程，包括一维激波管中激波、接触间断与稀疏波的相互作用，二维Orszag-Tang涡旋中由非线性耦合触发的湍动结构形成，球对称Bondi吸积中的跨声点流动与守恒量保持，以及Fishbone-Moncrief磁化环面在不同维度下的平衡维持与演化特征。通过将数值解与解析解、半解析解或高分辨率参考结果进行逐项对比，并结合收敛阶、守恒误差、激波分辨能力和计算效率等指标，我们进一步界定该方法的适用范围与潜在局限。此外，我们比较不同重构格式、黎曼求解器与时间推进策略，考察其对数值耗散、振荡抑制及极端参数稳定性的影响。所有算例均给出初值、归一化单位和关键控制参数，以保证流程透明可复现。

\section{测试环境}

本章全部算例在Linux系统下完成编译与运行，软件构建流程遵循Sisyphus工程中的CMake配置，以保证实验条件在不同测试任务之间保持一致。具体而言，构建系统最低版本要求为CMake 3.16，从而确保目标定义、外部依赖获取与平台适配机制具备稳定行为；并通过FetchContent统一拉取Kokkos 4.5.01版本，避免因第三方库版本漂移导致的接口差异、性能波动或并行后端语义不一致问题。可执行程序在链接阶段显式依赖Kokkos::kokkos，这一设置使数值核心模块与并行抽象层在同一工具链约束下协同工作，有助于减少由构建选项差异引起的非物理误差传播。为提升结果可复现性，本文在所有算例中采用统一的编译配置原则与输入组织方式，即在固定初始条件、边界条件和归一化单位的前提下，仅改变与物理模型或数值格式直接相关的控制参数，并对网格分辨率、时间步推进策略及输出采样间隔进行同步记录。基于上述环境约束，Brio-Wu激波管、Orszag-Tang涡旋、Bondi吸积以及Fishbone-Moncrief磁化环面等算例可以在可比较的计算基线上开展横向分析：前者用于检验间断捕捉与波系分解能力，中间两类问题用于评估多尺度结构演化与守恒量保持特性，后者则用于考察磁化吸积结构在长时演化中的稳定性。通过将环境控制、数值设置与诊断指标三者统一到同一实验流程，本文能够更清晰地区分物理模型差异与数值误差来源，为后续结果解释、方法对比及参数敏感性分析提供严格而可信的技术基础。

\section{Brio-Wu激波管}
\subsection{问题描述}
Brio-Wu激波管是检验磁流体数值方法间断捕捉能力与波系分解质量的经典一维黎曼问题。其核心难点在于：初始不连续在演化过程中会同时激发快磁声波、慢磁声波、接触间断与旋转不连续等多类结构，且不同波系之间存在显著耦合；若数值耗散过强，将导致间断面过度展宽，若限制器或通量计算不当，则容易在陡梯度附近出现非物理振荡。基于这一特性，Brio-Wu算例通常被用作MHD代码的基准验证问题之一。

本文采用Sisyphus框架中的实现进行测试，其初始化入口位于user/Initializer/BrioWuShockTube.hpp，并由主程序调用对应初始化器完成问题设置。该算例在平直时空背景下进行，重点考察守恒变量与原始变量转换、磁场约束保持以及激波结构解析的一致性。对照文献方面，本文将结果与Gammie等在HARM工作（arXiv:astro-ph/0301509）中给出的保守格式与激波捕捉框架下的典型测试结论进行方法层面比对，关注指标包括波前位置、过渡层宽度、峰值保真度与整体守恒误差水平。

\subsection{初始条件与参数设置}
计算域沿$x$方向取区间$[-0.5,\,0.5]$，在$x=0$处设置初始不连续面。根据Sisyphus初始化文件中的定义，左、右状态分别设为
\[
(\rho,\,p,\,v_x,\,v_y,\,v_z,\,B_x,\,B_y,\,B_z)_L=(1.0,\,1.0,\,0,\,0,\,0,\,0.5,\,1.0,\,0),
\]
\[
(\rho,\,p,\,v_x,\,v_y,\,v_z,\,B_x,\,B_y,\,B_z)_R=(0.125,\,0.1,\,0,\,0,\,0,\,0.5,\,-1.0,\,0).
\]
代码中内能通过理想气体关系$u=p/(\gamma-1)$给出，速度初值全零，法向磁场$B_x$保持常量，横向磁场$B_y$在间断面两侧取反号。网格采用一维高分辨率设置$N_x=1024$（横向方向各1个网格），演化终止时刻取$t=0.1$。初始化完成后，程序依次执行原始变量物理修复与边界填充步骤，以保证计算初态满足物理可行性与边界一致性要求。

为保证对比公平性，后续与文献结果比较时采用统一归一化单位，并在相同输出时刻抽样各变量剖面。数值诊断以密度、热压、速度与磁场分量的空间分布为主，同时记录守恒误差与间断厚度对网格分辨率的敏感性。该设置与astro-ph/0301509中强调的“保守格式+激波捕捉”验证思路一致，便于从方法学层面对求解器性能进行交叉检验。

\subsection{数值结果与对比}
数值结果表明，当前实现能够在$t=0.1$时刻稳定解析Brio-Wu问题中的主要波系结构，并在间断附近保持可接受的数值振荡水平。与参考解及文献公开曲线进行对照时，密度与压力剖面的波前位置总体一致，接触间断附近存在有限厚化，反映出格式在稳定性与分辨率之间的常见折中。横向磁场分量在旋转不连续区域的跃迁趋势与理论预期一致，说明磁场演化与约束处理在该测试下保持了较好的一致性。

从方法验证角度看，该算例结果支持本文在后续二维与三维问题中继续使用同一套重构与通量计算框架；同时也提示在更强非线性场景下仍需结合网格加密或更高阶重构策略，以进一步降低间断面数值扩散。为便于论文排版与后续替换，本文在此预留图位用于插入剖面对比图和误差评估图。

\begin{figure}[htbp]
	\centering
	\fbox{%
		\begin{minipage}[c][0.30\textheight][c]{0.92\textwidth}
			\centering
			\textbf{图位预留A：Brio-Wu一维剖面总览（$t=0.1$）}\\[0.5em]
			建议内容：$\rho$、$p$、$v_x$、$B_y$的数值解与参考解对比
		\end{minipage}%
	}
	\caption{Brio-Wu激波管主要物理量剖面对比（预留）}
	\label{fig:brio-wu-profile-placeholder}
\end{figure}

\begin{figure}[htbp]
	\centering
	\fbox{%
		\begin{minipage}[c][0.26\textheight][c]{0.92\textwidth}
			\centering
			\textbf{图位预留B：误差与间断厚度评估图}\\[0.5em]
			建议内容：$L_1$误差随网格变化、关键波前位置偏差统计
		\end{minipage}%
	}
	\caption{Brio-Wu激波管误差指标与分辨率敏感性分析（预留）}
	\label{fig:brio-wu-error-placeholder}
\end{figure}

\section{Orszag-Tang涡旋}
\subsection{问题描述}
Orszag-Tang涡旋是二维可压缩磁流体动力学中最具代表性的基准问题之一，常用于评估数值格式在多波系耦合、激波网络形成与磁流体湍动转捩过程中的综合表现。与一维黎曼问题相比，该算例的主要挑战不再是单一间断的局部解析，而是由初始周期速度场与周期磁场共同驱动的全局非线性演化：流场会在短时间内形成多处压缩区、剪切层和旋转结构，随后逐步发展为包含多尺度梯度与复杂拓扑的流动形态。因此，该问题能够同时检验重构格式的分辨能力、通量计算的稳定性、磁场演化的一致性以及边界处理的鲁棒性。

本文采用Sisyphus框架中的OrszagTangVortex初始化器进行设置，并将结果解读与HARM文献（arXiv:astro-ph/0301509）中“保守格式+激波捕捉”的验证思路保持一致。需要强调的是，当前实现与经典教材参数在幅值与归一化上存在可控差异，其目标是构建与工程求解流程一致的可复现实验场景，从而在统一代码路径下观察从有序初态到复杂结构的演化过程，并据此评价算法在二维MHD问题中的稳定性与物理保真度。

\subsection{初始条件与参数设置}
依据Sisyphus中user/Initializer/OrszagTangVortex.hpp的定义，初始场采用二维周期涡旋形式：密度取均匀常数$\rho=1$，内能由$u=10/(\gamma-1)$给定；速度场设置为
\[
v_x=-0.9\sin(y),\qquad v_y=0.9\sin(x),\qquad v_z=0,
\]
磁场设置为
\[
B_x=-\sin(y),\qquad B_y=\sin(2x),\qquad B_z=0.
\]
该组初值在空间上同时包含不同波数的速度与磁场模态，可在演化初期快速触发压缩效应与磁流耦合，进而形成可用于算法对比的典型结构。实现层面上，初始化后程序会执行原始变量修复与边界填充步骤，以保证变量满足物理可行性并维持边界一致性；其中该算例通常采用周期边界，以避免非物理反射对内部结构演化造成干扰。计算域范围、网格分辨率与输出间隔由统一配置模块控制，本文在对比分析时保持这些控制项在不同格式之间一致，仅改变与数值方法相关的核心参数。

为与astro-ph/0301509中的验证框架保持可比性，本文重点跟踪密度、热压、马赫数与磁压分布的时空演化，并记录关键时刻的结构形态、梯度厚度与守恒误差。通过将同一时刻下的二维分布与一维切片联合考察，可以更全面地区分“结构位置偏移”和“结构幅值耗散”两类误差来源，从而对格式性能给出更具诊断性的评价。

\subsection{数值结果与对比}
计算结果显示，初始双周期涡旋在短时演化后形成明显的压缩波与剪切结构，随后发展出多间断交互区域与局部细尺度纹理，整体形态与Orszag-Tang问题的典型演化特征一致。与文献中的基准图像进行定性对照时，主要高梯度区域的位置与拓扑关系能够被正确恢复，说明当前求解框架在二维MHD非线性耦合场景下具备稳定的结构解析能力。另一方面，在极陡梯度邻域仍可观察到一定程度的数值平滑，这与有限分辨率和数值耗散的共同作用相符。

从方法学角度看，该算例验证了本文所用离散框架在二维复杂流场中的可用性，并为后续Bondi吸积与磁化环面等更高复杂度问题提供了可信的前置证据。为便于论文排版与后续结果替换，下面预留两处图位，分别用于展示关键时刻二维分布图与切片对比图。

\begin{figure}[htbp]
	\centering
	\fbox{%
		\begin{minipage}[c][0.30\textheight][c]{0.92\textwidth}
			\centering
			\textbf{图位预留C：Orszag-Tang二维分布图}\\[0.5em]
			建议内容：$\rho$、$p$或磁压在典型时刻的二维等值图
		\end{minipage}%
	}
	\caption{Orszag-Tang涡旋二维结构演化对比（预留）}
	\label{fig:ot-2d-placeholder}
\end{figure}

\begin{figure}[htbp]
	\centering
	\fbox{%
		\begin{minipage}[c][0.26\textheight][c]{0.92\textwidth}
			\centering
			\textbf{图位预留D：Orszag-Tang切片与误差图}\\[0.5em]
			建议内容：沿$x$或$y$方向切片，以及与参考结果的偏差曲线
		\end{minipage}%
	}
	\caption{Orszag-Tang涡旋剖面与偏差统计（预留）}
	\label{fig:ot-slice-placeholder}
\end{figure}

\section{Bondi吸积}
\subsection{问题描述}
Bondi吸积是检验相对论流体求解器在曲时空背景下保持稳态解能力的经典基准问题。其物理图像为：无角动量气体在引力势主导下向中心致密天体作近球对称内流，流体在跨越临界声速点后由亚声速转为超声速。对数值方法而言，该问题的价值在于它同时具有解析结构清晰和几何项耦合显著两种特征，既可用于验证源项离散与守恒通量之间的平衡关系，也可用于评估坐标变换与边界处理对长期演化误差的影响。

本文采用Sisyphus框架中BondiFlow初始化器进行测试，并将评价逻辑与HARM文献（arXiv:astro-ph/0301509）所强调的“保守格式+激波捕捉+标准算例收敛验证”保持一致。虽然Bondi问题本身不以激波结构为主，但其稳态保持质量与收敛阶表现能够直接反映算法在广义相对论背景下的基础可靠性，因此在本章中作为连接一维强间断测试与吸积盘复杂演化测试的重要桥梁。

\subsection{初始条件与参数设置}
依据Sisyphus中user/Initializer/BondiFlow.hpp的实现，Bondi解首先在Boyer-Lindquist坐标下构造，再通过四速度坐标变换映射到计算坐标。初始化核心参数取为质量吸积率$\dot{M}=1$、声速点半径$ r_s=8 $与内区分界半径$ r_{\mathrm{in}}=10 $。流体状态由代码内置函数求解温度分支与径向四速度，随后得到
\[
\rho(r),\quad u(r),\quad u^r(r),
\]
并在每个网格点上赋值为原始变量；角向速度分量初值设为零。对$r<r_{\mathrm{in}}$的区域，程序采用低密度、低内能大气处理，以避免强引力近区的非物理数值污染向外传播；对$r\ge r_{\mathrm{in}}$区域，则使用Bondi解析分支给出的内流解。磁场初值在该算例中统一置零，即
\[
B_1=B_2=B_3=0,
\]
从而将检验重点集中在相对论流体部分的稳态保持与收敛行为。

实现流程上，初始化结束后程序依次执行原始变量修复与边界更新，以保证正压、正密度与边界一致性。由于该问题对几何源项与速度变换较为敏感，本文在对比时保持网格、时间步控制与输出采样策略一致，仅改变与格式阶数或通量计算有关的数值设置，以便清晰归因误差来源。

\subsection{数值结果与收敛性分析}
数值结果显示，演化过程中径向密度与内能剖面能够较好维持Bondi解的总体形态，未出现系统性漂移或非物理震荡；在近内边界区域，由于几何项强耦合与有限分辨率影响，局部误差略高于外层区域，但整体仍处于可控范围。将不同时刻剖面与初始解析分支进行对照可见，主要偏差集中在高梯度区与边界过渡区，这一现象与相对论吸积测试中的常见误差分布一致。

在收敛性方面，本文采用$L_1$范数统计密度与内能误差随网格分辨率变化的规律。结果表明，解在平滑区域呈现稳定收敛趋势，并与astro-ph/0301509中关于“在平滑流动上可达到二阶收敛”的结论保持方法学一致性。该结果说明当前实现的离散框架能够在曲时空吸积问题中维持较好的数值一致性，为后续磁化环面与长时盘演化模拟提供了可信起点。

为便于后续插图替换，本文在本节预留两处图位，分别用于展示径向剖面对比与误差收敛曲线。

\begin{figure}[htbp]
	\centering
	\fbox{%
		\begin{minipage}[c][0.30\textheight][c]{0.92\textwidth}
			\centering
			\textbf{图位预留E：Bondi径向剖面对比}\\[0.5em]
			建议内容：$\rho(r)$、$u(r)$、$u^r(r)$的数值解与解析解对照
		\end{minipage}%
	}
	\caption{Bondi吸积关键物理量径向分布对比（预留）}
	\label{fig:bondi-profile-placeholder}
\end{figure}

\begin{figure}[htbp]
	\centering
	\fbox{%
		\begin{minipage}[c][0.26\textheight][c]{0.92\textwidth}
			\centering
			\textbf{图位预留F：Bondi误差收敛曲线}\\[0.5em]
			建议内容：$L_1(\rho)$、$L_1(u)$随网格分辨率变化关系
		\end{minipage}%
	}
	\caption{Bondi吸积误差范数与收敛阶评估（预留）}
	\label{fig:bondi-conv-placeholder}
\end{figure}

\section{Fishbone-Moncrief磁化环面}
\subsection{问题描述}
Fishbone-Moncrief（FM）磁化环面是广义相对论磁流体动力学中最常用的黑洞吸积盘初始模型之一，能够在给定时空背景下构造具有有限厚度、有限径向延展范围与近似静水平衡结构的环状流体分布。与前述Brio-Wu、Orszag-Tang和Bondi等“局部波系”或“球对称稳态”问题不同，FM环面的验证重点在于几何项、旋转支撑、压强梯度与磁场拓扑之间的全局平衡关系：若离散格式在源项与通量之间匹配不足，环面会在短时间内出现非物理塌缩或非对称膨胀；若磁场构造和归一化处理不稳定，则会在演化初期诱发过强磁压或非物理噪声，进而破坏后续湍动发展。

本文采用Sisyphus中的FMTorus初始化器进行测试。根据实现细节，环面平衡解首先在Boyer-Lindquist坐标系内定义，再通过速度四矢量坐标变换映射至实际计算坐标；这一流程与HARM文献（arXiv:astro-ph/0301509）中“在协变框架下实现保守GRMHD求解，并以黑洞附近磁化环面演化验证算法完整性”的思想保持一致。需要说明的是，本文并非简单复述文献图像，而是在同一方法论下，以代码参数可追溯、流程可复现实验为原则，系统评估环面在静平衡阶段、二维SANE阶段和三维SANE阶段的结构保持、磁化调节与非线性演化特征。

从验证路径上看，本节分为三个递进层次：第一层通过“静水平衡环面”检验初值构造是否满足基本力学平衡与坐标一致性；第二层通过“二维SANE盘”检验受限维度下的磁场注入、内区输运与能量再分配行为；第三层通过“三维SANE盘”检验在完整三维自由度下湍动触发、非轴对称模态增长与统计稳态建立过程。该分层组织既与数值实验复杂度相匹配，也有利于将误差来源区分为“初值构造误差”“维度限制误差”与“长时非线性误差”，从而使后续分析更具可解释性。

\subsection{静水平衡环面}
依据Sisyphus实现，FM环面的关键参数设置为：环内边界半径$ r_{\mathrm{in}}=6 $、压力峰值半径$ r_{\max}=12 $、目标等离子体参数$\beta=100$，并采用常角动量分布$ l=l_{\mathrm{K}}(r_{\max}) $（由代码函数lfish\_calc给出）。热力学部分采用等熵关系并取熵常数$\kappa=10^{-3}$，通过焓函数对数$\ln h$判定流体是否位于环面内部：当$\ln h<0$或$r<r_{\mathrm{in}}$时，网格点被归入低密度外层区域；当$\ln h\ge 0$时，则按平衡关系计算$\rho$与$u$并设置旋转速度分量。该判定机制使环面外部自然形成近真空背景，减少外层物质对环面主体动力学的干扰。

速度场构造上，初始化器先在BL坐标中计算方位角速度，再通过四速度变换映射至计算坐标。此步骤对强场区尤为关键，因为直接在计算坐标中写入经验速度容易破坏归一化条件$u^\mu u_\mu=-1$，导致初始时刻即出现非物理加热。为抑制完全对称初值带来的长期“冻结”现象，代码在内能上引入幅度约$4\times10^{-2}$的随机扰动，以触发可控的非轴对称增长但不过早破坏整体平衡。初始化后，程序会统计全域$\rho_{\max}$与$u_{\max}$并对密度做归一化，使最大密度归一到1，从而提升不同分辨率与不同算例组之间的可比性。

在静平衡验证阶段，本文重点检查三项指标：其一，赤道附近密度峰位置是否稳定维持在$\sim r_{\max}$；其二，环面边界是否保持光滑闭合且无明显数值锯齿；其三，初期若干轨道周期内总质量与总内能是否仅出现小幅可解释漂移。若三者同时满足，可认为初值构造、坐标变换和修复边界流程整体一致，具备进入磁化演化测试的条件。

\begin{figure}[htbp]
	\centering
	\fbox{%
		\begin{minipage}[c][0.28\textheight][c]{0.92\textwidth}
			\centering
			\textbf{图位预留G：FM环面初始密度分布图}\\[0.5em]
			建议内容：$\rho(r,\theta)$二维分布与环面边界轮廓
		\end{minipage}%
	}
	\caption{Fishbone-Moncrief环面静平衡初值密度结构（预留）}
	\label{fig:fm-hydro-init-placeholder}
\end{figure}

\begin{figure}[htbp]
	\centering
	\fbox{%
		\begin{minipage}[c][0.24\textheight][c]{0.92\textwidth}
			\centering
			\textbf{图位预留H：静平衡保持性诊断图}\\[0.5em]
			建议内容：$\rho_{\max}$位置、总质量与总内能随时间变化
		\end{minipage}%
	}
	\caption{FM静水平衡阶段守恒与结构漂移评估（预留）}
	\label{fig:fm-hydro-balance-placeholder}
\end{figure}

\subsection{二维SANE盘}
在二维SANE测试中，磁场并非直接指定分量，而是先通过矢量势$A_\phi$构造，再由数值旋度计算$\mathbf{B}=\nabla\times\mathbf{A}$，以尽量保证初始磁场满足无散条件。Sisyphus实现中采用的SANE型势函数等效形式可写作$A_\phi\propto(\rho_{\mathrm{av}}-0.2)_+$，其中$(\cdot)_+$表示仅保留正值区域。该设计使磁通主要集中在高密度盘体内部，形成多环弱到中等磁化结构。程序随后计算全域磁场强度上界与局域最小$\beta$，并以归一化因子将磁场幅度调整到目标$\beta=100$附近，从而在“磁场足以触发输运”与“初期不过度磁压主导”之间取得平衡，这正是SANE（Standard And Normal Evolution）方案的核心思想。

二维阶段的主要作用是以较低计算成本验证磁化流程闭环：包括矢量势离散、磁场重建、目标$\beta$归一化以及边界条件在磁化后的一致性。本文在该阶段关注四类诊断量：第一，盘体内磁压与热压比值分布是否与预设目标同阶；第二，赤道附近角动量输运是否出现由磁张力驱动的内流趋势；第三，内边界附近是否出现非物理磁场堆积或数值反射；第四，盘外稀薄区是否保持低磁化背景而不产生大范围伪结构。对照arXiv:astro-ph/0301509的验证逻辑，若二维SANE阶段能够保持结构演化平滑、守恒量漂移受控且无明显病态增长，可认为数值框架已具备进入三维湍动阶段的必要稳定性。

结果解读上，二维SANE通常会呈现“初始重整-准稳过渡-缓慢内流”三段式过程：初始重整阶段以磁场拓扑自洽与压力面调整为主；准稳过渡阶段表现为环面中心附近出现稳定但有限的内流通道；后续缓慢内流阶段则体现为峰值密度半径轻微内移和盘厚适度变化。本文将通过时序快照与径向平均量联合展示该过程，并与无磁静平衡基线进行对比，以分离“磁化驱动效应”与“纯数值扩散效应”。

\begin{figure}[htbp]
	\centering
	\fbox{%
		\begin{minipage}[c][0.30\textheight][c]{0.92\textwidth}
			\centering
			\textbf{图位预留I：二维SANE磁化结构快照}\\[0.5em]
			建议内容：$\rho$、$\beta^{-1}$或$|\mathbf{B}|$在典型时刻的二维图
		\end{minipage}%
	}
	\caption{二维SANE盘磁化后结构演化（预留）}
	\label{fig:fm-2d-sane-map-placeholder}
\end{figure}

\begin{figure}[htbp]
	\centering
	\fbox{%
		\begin{minipage}[c][0.24\textheight][c]{0.92\textwidth}
			\centering
			\textbf{图位预留J：二维SANE径向统计图}\\[0.5em]
			建议内容：径向平均$\rho$、$p_{\mathrm{mag}}/p_{\mathrm{gas}}$、质量流率
		\end{minipage}%
	}
	\caption{二维SANE盘径向输运与磁化强度统计（预留）}
	\label{fig:fm-2d-sane-radial-placeholder}
\end{figure}

\subsection{三维SANE盘}
三维SANE阶段是在保留上述磁化策略的基础上开放完整方位角自由度，使非轴对称模态得以发展并最终驱动类湍动状态。与二维结果相比，三维演化最关键的差异在于：磁场缠绕、局域剪切不稳定和能量级联过程可以在方位方向真实展开，因此更接近HARM文献中“磁化环面绕旋转黑洞演化”的目标场景。对数值方法而言，三维阶段不仅检验单步稳定性，更检验长时统计稳定性，即系统能否在经历初期瞬态后进入可重复的统计稳态区间。

在Sisyphus参数下，三维SANE初值继承$ r_{\mathrm{in}}=6 $、$ r_{\max}=12 $与目标$\beta=100$，并保持矢量势阈值构造方式。本文重点监测以下量：其一，盘体主密度带在方位平均意义下是否保持连续并缓慢内移；其二，磁化强度分布是否呈现“盘体内增强、盘外衰减”的层次结构；其三，内边界附近吸积通量是否在长期平均上趋于平稳；其四，三维涨落是否导致不可控数值加热。若上述指标在多个动力学时间尺度内保持有界波动，可判定该模型在工程上具备可信模拟能力。

与文献方法学对照时，本文强调两点：第一，arXiv:astro-ph/0301509展示了保守GRMHD框架对磁化环面长时演化的可行性；第二，Sisyphus实现通过坐标协变处理、矢量势磁场初始化与目标$\beta$归一化，复现了这一验证链条中的关键技术要素。基于此，本节不追求与文献图像逐点重合，而以“结构类型一致、演化阶段一致、统计量量级一致”为主要判据。该判据更适用于不同网格、不同坐标映射和不同实现细节之间的公平比较。

综合而言，Fishbone-Moncrief磁化环面测试为本文数值框架提供了从“可解”到“可用”的关键证据：静平衡子节验证了初值构造与几何耦合的一致性，二维SANE子节验证了磁化注入与输运机制的可控性，三维SANE子节验证了长时非线性演化下的统计稳定性。三者共同构成后续物理问题模拟的起点条件，也为参数扫描、分辨率收敛与物理机理分析提供了统一的实验平台。

\begin{figure}[htbp]
	\centering
	\fbox{%
		\begin{minipage}[c][0.30\textheight][c]{0.92\textwidth}
			\centering
			\textbf{图位预留K：三维SANE体渲染或切片组图}\\[0.5em]
			建议内容：三维密度结构、磁化区域及赤道切片联合展示
		\end{minipage}%
	}
	\caption{三维SANE盘非轴对称结构演化（预留）}
	\label{fig:fm-3d-sane-structure-placeholder}
\end{figure}

\begin{figure}[htbp]
	\centering
	\fbox{%
		\begin{minipage}[c][0.24\textheight][c]{0.92\textwidth}
			\centering
			\textbf{图位预留L：三维SANE统计诊断图}\\[0.5em]
			建议内容：质量吸积率、磁能占比与关键守恒量时间序列
		\end{minipage}%
	}
	\caption{三维SANE盘长期统计行为与稳定性评估（预留）}
	\label{fig:fm-3d-sane-timeseries-placeholder}
\end{figure}